\documentclass[11pt]{article}
\usepackage{amsmath,amssymb,amsthm}
\usepackage{algorithm,algorithmic}
\usepackage{enumitem}
\usepackage{hyperref}
\usepackage{geometry}
\geometry{margin=1in}
\newtheorem{theorem}{Theorem}
\newtheorem{lemma}{Lemma}
\newtheorem{assumption}{Assumption}
\newtheorem{remark}{Remark}
\newcommand{\E}{\mathbb{E}}
\newcommand{\R}{\mathbb{R}}

\begin{document}

\title{Convergence Analysis of a Scalar AdaGrad‑Style Gradient–Average Method}
\author{}
\date{}
\maketitle

%%%%%%%%%%%%%%%%%%%%%%%%%%%%%%%%%%%%%%%%%%%%%%%%%%%%%%%%%%%%%%%%%
\section*{Algorithm Name}
\textbf{Scalar AdaGrad with Gradient Accumulator}\\
(also referred to as the \emph{running‑average AdaGrad} method)

%%%%%%%%%%%%%%%%%%%%%%%%%%%%%%%%%%%%%%%%%%%%%%%%%%%%%%%%%%%%%%%%%
\section*{Mathematical Formulation}
Let $f:\R\rightarrow\R$ be a convex, $L$‑smooth function.  
Denote by $g_k=\nabla f(w_k)$ the (deterministic) gradient evaluated at iteration $k$.
Define the cumulative gradient sum and the cumulative squared‑gradient sum as
\[
A_k \;:=\; \sum_{i=1}^{k} g_i ,\qquad
G_k \;:=\; \sum_{i=1}^{k} g_i^{2}.
\]
Given a positive stepsize hyperparameter $\alpha>0$, the update rule is
\begin{equation}
\boxed{
w_{k+1}= w_k-\frac{\alpha}{\sqrt{G_k+\varepsilon}}\; g_k},
\qquad k=0,1,2,\dots,
\label{eq:update}
\end{equation}
where $\varepsilon>0$ is a small constant (e.g. $10^{-8}$) introduced for numerical stability.  
The denominator $\sqrt{G_k+\varepsilon}$ is the element‑wise AdaGrad scaling; in the scalar case it reduces to the square root of the accumulated squared gradients.

%%%%%%%%%%%%%%%%%%%%%%%%%%%%%%%%%%%%%%%%%%%%%%%%%%%%%%%%%%%%%%%%%
\section*{Pseudocode}
\begin{algorithm}[H]
\caption{Scalar AdaGrad with Gradient Accumulator}
\begin{algorithmic}[1]
\REQUIRE Initial point $w_0\in\R$, stepsize $\alpha>0$, stability constant $\varepsilon>0$
\STATE $G_0 \gets 0$ \COMMENT{accumulated squared‑gradient}
\FOR{$k=0,1,\dots,T-1$}
    \STATE Compute gradient $g_k \gets \nabla f(w_k)$
    \STATE Update accumulator $G_{k+1} \gets G_k + g_k^{2}$
    \STATE Set adaptive stepsize $\eta_k \gets \dfrac{\alpha}{\sqrt{G_{k+1}+\varepsilon}}$
    \STATE Update iterate $w_{k+1} \gets w_k - \eta_k\, g_k$
\ENDFOR
\RETURN $w_T$
\end{algorithmic}
\end{algorithm}

%%%%%%%%%%%%%%%%%%%%%%%%%%%%%%%%%%%%%%%%%%%%%%%%%%%%%%%%%%%%%%%%%
\section*{Dry Run Example}
Consider the quadratic $f(w)=(w-3)^2$.  
Then $\nabla f(w)=2(w-3)$.  Choose $w_0=0$, $\alpha=0.1$, $\varepsilon=10^{-8}$ and run three iterations.

\begin{center}
\begin{tabular}{c|c|c|c|c}
$k$ & $w_k$ & $g_k=2(w_k-3)$ & $G_{k+1}=G_k+g_k^2$ & $w_{k+1}=w_k-\dfrac{\alpha}{\sqrt{G_{k+1}+\varepsilon}}g_k$\\\hline
0 & $0$ & $2(0-3)=-6$ & $0+(-6)^2=36$ &
$0-\dfrac{0.1}{\sqrt{36}}\cdot(-6)=0+\dfrac{0.1}{6}\cdot6=0.1$\\[4pt]
1 & $0.1$ & $2(0.1-3)=-5.8$ & $36+(-5.8)^2=36+33.64=69.64$ &
$0.1-\dfrac{0.1}{\sqrt{69.64}}\cdot(-5.8)
=0.1+\dfrac{0.1}{8.345}\cdot5.8\approx0.1695$\\[4pt]
2 & $0.1695$ & $2(0.1695-3)=-5.661$ & $69.64+(-5.661)^2\approx69.64+32.04=101.68$ &
$0.1695-\dfrac{0.1}{\sqrt{101.68}}\cdot(-5.661)
=0.1695+\dfrac{0.1}{10.084}\cdot5.661\approx0.2327$
\end{tabular}
\end{center}
Objective values:
\[
f(w_0)=9,\; f(w_1)= (0.1-3)^2\approx8.41,\; f(w_2)\approx8.27,\; f(w_3)\approx8.14.
\]
The iterates move toward the minimizer $w^\star=3$, with decreasing step sizes because $G_k$ grows.

%%%%%%%%%%%%%%%%%%%%%%%%%%%%%%%%%%%%%%%%%%%%%%%%%%%%%%%%%%%%%%%%%
\section*{Assumptions}
\begin{assumption}[Convexity]\label{as:convex}
$f:\R\rightarrow\R$ is convex, i.e.
\[
f(u)\ge f(v)+\nabla f(v)(u-v),\qquad\forall u,v\in\R .
\]
\end{assumption}
\begin{assumption}[Smoothness]\label{as:smooth}
$f$ has $L$‑Lipschitz continuous gradient:
\[
\bigl|\nabla f(u)-\nabla f(v)\bigr|\le L|u-v|,\qquad\forall u,v\in\R .
\]
\end{assumption}
\begin{assumption}[Bounded Domain]\label{as:domain}
All iterates lie in a closed interval $\mathcal{W}=[w_{\min},w_{\max}]$ of diameter
\[
D:= w_{\max}-w_{\min}<\infty .
\]
\end{assumption}
\begin{assumption}[Gradient Boundedness]\label{as:gradbound}
For all $w\in\mathcal{W}$,
\[
|\,\nabla f(w)\,|\le G_{\max } .
\]
\end{assumption}
All constants $L,G_{\max},D,\alpha,\varepsilon$ are deterministic and known.

%%%%%%%%%%%%%%%%%%%%%%%%%%%%%%%%%%%%%%%%%%%%%%%%%%%%%%%%%%%%%%%%%
\section*{Main Convergence Theorem}
\begin{theorem}[Regret / Optimization bound for scalar AdaGrad]\label{thm:main}
Under Assumptions \ref{as:convex}–\ref{as:gradbound}, let $w^\star\in\arg\min_{w\in\mathcal{W}} f(w)$.  
For the iterates $\{w_k\}_{k=0}^{T}$ generated by \eqref{eq:update} with any $\alpha>0$ and $\varepsilon>0$,
\[
\frac{1}{T}\sum_{k=0}^{T-1}\bigl(f(w_k)-f(w^\star)\bigr)
\;\le\;
\frac{D^{2}}{2\alpha}\,\frac{\sqrt{G_T+\varepsilon}}{T}
\;+\;
\frac{\alpha G_{\max }^{2}}{2}\,\frac{1}{T}\sum_{k=0}^{T-1}\frac{1}{\sqrt{G_k+\varepsilon}} .
\]
In particular, choosing $\alpha=D/\bigl(G_{\max }\sqrt{T}\bigr)$ yields
\[
\frac{1}{T}\sum_{k=0}^{T-1}\bigl(f(w_k)-f(w^\star)\bigr)
\;\le\;
\frac{DG_{\max }}{\sqrt{T}}+O\!\left(\frac{1}{T}\right),
\]
i.e. an $O(1/\sqrt{T})$ convergence rate.
\end{theorem}

%%%%%%%%%%%%%%%%%%%%%%%%%%%%%%%%%%%%%%%%%%%%%%%%%%%%%%%%%%%%%%%%%
\section*{Theoretical Properties}
\begin{enumerate}[label=\arabic*.]
\item \textbf{Stability.}  
Because the denominator $\sqrt{G_k+\varepsilon}$ is non‑decreasing, the adaptive stepsize $\eta_k$ never exceeds $\alpha/\sqrt{\varepsilon}$, guaranteeing bounded updates.

\item \textbf{Hyper‑parameter sensitivity.}  
$\alpha$ scales the whole bound linearly; a larger $\alpha$ reduces the first term (initial distance) but inflates the second term (gradient variance). The optimal $\alpha$ balances the two, leading to the $1/\sqrt{T}$ rate.

\item \textbf{Comparison to SGD with fixed stepsize.}  
For SGD with a constant $\eta$, the standard bound is $O(1/\sqrt{T})$ only under a diminishing stepsize schedule. AdaGrad automatically adapts the stepsize based on observed gradients, achieving the same order without an explicit schedule.

\item \textbf{Effect of smoothness.}  
If $f$ is $L$‑smooth, one can replace the bounded‑gradient assumption by $|g_k|\le L D$, yielding identical constants with $G_{\max}=L D$.

\end{enumerate}

%%%%%%%%%%%%%%%%%%%%%%%%%%%%%%%%%%%%%%%%%%%%%%%%%%%%%%%%%%%%%%%%%
\section*{Discussion}
The theorem shows that the scalar AdaGrad method enjoys the canonical $O(1/\sqrt{T})$ convergence for convex smooth objectives.  
The proof technique mirrors the classical AdaGrad analysis for vector‑valued parameters, but the scalar setting allows a fully explicit derivation.  
The adaptive denominator yields smaller steps as the accumulated squared gradient grows, which is precisely what stabilises the method on ill‑conditioned problems.

%%%%%%%%%%%%%%%%%%%%%%%%%%%%%%%%%%%%%%%%%%%%%%%%%%%%%%%%%%%%%%%%%
\section*{Preliminaries (Definitions and Lemmas)}
\begin{definition}[Regret]
For a sequence $\{w_k\}_{k=0}^{T-1}$ and a comparator $w^\star\in\mathcal{W}$,
\[
\text{Regret}_T:=\sum_{k=0}^{T-1}\bigl(f(w_k)-f(w^\star)\bigr).
\]
\end{definition}
\begin{lemma}[Standard convexity inequality]\label{lem:conv}
For any convex $f$ and any $w,w^\star$,
\[
f(w)-f(w^\star)\le \nabla f(w)\,(w-w^\star).
\]
\end{lemma}
\begin{lemma}[Quadratic bound on distance]\label{lem:dist}
For any $a,b\in\R$ and any $\eta>0$,
\[
(a-b)^2 = (a-c)^2 -2\eta (a-c)(c-b)+\eta^2(b-c)^2
\]
holds with $c:=a-\eta b$.
\end{lemma}
\begin{proof}
Direct expansion of $(a-\eta b-b)^2$ yields the identity. \hfill $\square$
\end{proof}

%%%%%%%%%%%%%%%%%%%%%%%%%%%%%%%%%%%%%%%%%%%%%%%%%%%%%%%%%%%%%%%%%
\section*{Main Proof (Step‑by‑Step Derivation)}
We prove Theorem~\ref{thm:main}.  All steps are numbered for easy reference.

\bigskip
\noindent\textbf{Step 1 (Convexity).}  
From Lemma~\ref{lem:conv} applied with $w=w_k$,
\begin{equation}
f(w_k)-f(w^\star) \le g_k\,(w_k-w^\star). \tag{1}
\end{equation}

\noindent\textbf{Step 2 (Expand the distance after the update).}  
Define $\eta_k:=\alpha/\sqrt{G_{k+1}+\varepsilon}$ and recall the update
$w_{k+1}=w_k-\eta_k g_k$.  
Using Lemma~\ref{lem:dist} with $a=w_k$, $b=g_k$, $c=w_{k+1}$ and $\eta=\eta_k$,
\begin{equation}
\|w_{k+1}-w^\star\|^{2}
= \|w_k-w^\star\|^{2}
-2\eta_k g_k (w_k-w^\star)
+\eta_k^{2} g_k^{2}. \tag{2}
\end{equation}

\noindent\textbf{Step 3 (Rearrange to isolate the inner product).}  
Move the middle term to the left side of (2):
\[
2\eta_k g_k (w_k-w^\star)
= \|w_k-w^\star\|^{2}
-\|w_{k+1}-w^\star\|^{2}
+\eta_k^{2} g_k^{2}. \tag{3}
\]

\noindent\textbf{Step 4 (Plug (3) into the convexity bound (1)).}  
Multiplying (1) by $2\eta_k$ and substituting the right‑hand side of (3) gives
\[
2\eta_k\bigl(f(w_k)-f(w^\star)\bigr)
\le
\|w_k-w^\star\|^{2}
-\|w_{k+1}-w^\star\|^{2}
+\eta_k^{2} g_k^{2}. \tag{4}
\]

\noindent\textbf{Step 5 (Divide by $2\eta_k$).}  
Since $\eta_k>0$,
\[
f(w_k)-f(w^\star)
\le
\frac{\|w_k-w^\star\|^{2}-\|w_{k+1}-w^\star\|^{2}}{2\eta_k}
+\frac{\eta_k}{2} g_k^{2}. \tag{5}
\]

\noindent\textbf{Step 6 (Insert the explicit form of $\eta_k$).}  
Recall $\eta_k = \dfrac{\alpha}{\sqrt{G_{k+1}+\varepsilon}}$.
Thus,
\[
\frac{1}{\eta_k}= \frac{\sqrt{G_{k+1}+\varepsilon}}{\alpha},\qquad
\eta_k = \frac{\alpha}{\sqrt{G_{k+1}+\varepsilon}}.
\]
Plugging into (5) yields
\[
f(w_k)-f(w^\star)
\le
\frac{\sqrt{G_{k+1}+\varepsilon}}{2\alpha}\bigl(\|w_k-w^\star\|^{2}-\|w_{k+1}-w^\star\|^{2}\bigr)
+\frac{\alpha}{2\sqrt{G_{k+1}+\varepsilon}}\, g_k^{2}. \tag{6}
\]

\noindent\textbf{Step 7 (Sum over $k=0,\dots,T-1$).}  
Summing (6) telescopically gives
\[
\sum_{k=0}^{T-1}\bigl(f(w_k)-f(w^\star)\bigr)
\le
\frac{1}{2\alpha}\sum_{k=0}^{T-1}\sqrt{G_{k+1}+\varepsilon}\,\bigl(\|w_k-w^\star\|^{2}-\|w_{k+1}-w^\star\|^{2}\bigr)
+\frac{\alpha}{2}\sum_{k=0}^{T-1}\frac{g_k^{2}}{\sqrt{G_{k+1}+\varepsilon}}. \tag{7}
\]

\noindent\textbf{Step 8 (Upper bound the telescoping term).}  
Since $\sqrt{G_{k+1}+\varepsilon}\ge\sqrt{G_{k}+\varepsilon}$, the coefficients are non‑decreasing.  
Therefore,
\[
\sqrt{G_{k+1}+\varepsilon}\,(\|w_k-w^\star\|^{2}-\|w_{k+1}-w^\star\|^{2})
\le
\sqrt{G_{T}+\varepsilon}\,(\|w_k-w^\star\|^{2}-\|w_{k+1}-w^\star\|^{2}).
\]
Insert this bound into the first sum of (7):
\[
\frac{1}{2\alpha}\sqrt{G_{T}+\varepsilon}\sum_{k=0}^{T-1}\bigl(\|w_k-w^\star\|^{2}-\|w_{k+1}-w^\star\|^{2}\bigr)
=
\frac{1}{2\alpha}\sqrt{G_{T}+\varepsilon}\,\bigl(\|w_0-w^\star\|^{2}-\|w_T-w^\star\|^{2}\bigr). \tag{8}
\]

\noindent\textbf{Step 9 (Bound the distances).}  
By Assumption~\ref{as:domain}, $\|w_k-w^\star\|\le D$ for all $k$, hence
\[
\|w_0-w^\star\|^{2}\le D^{2},\qquad
\|w_T-w^\star\|^{2}\ge 0.
\]
Thus (8) is bounded by $\dfrac{D^{2}}{2\alpha}\sqrt{G_{T}+\varepsilon}$.

\noindent\textbf{Step 10 (Bound the second sum).}  
Since $|g_k|\le G_{\max }$ (Assumption~\ref{as:gradbound}),
\[
\frac{g_k^{2}}{\sqrt{G_{k+1}+\varepsilon}}\le
\frac{G_{\max }^{2}}{\sqrt{G_{k+1}+\varepsilon}}.
\]
Hence
\[
\frac{\alpha}{2}\sum_{k=0}^{T-1}\frac{g_k^{2}}{\sqrt{G_{k+1}+\varepsilon}}
\le
\frac{\alpha G_{\max }^{2}}{2}\sum_{k=0}^{T-1}\frac{1}{\sqrt{G_{k+1}+\varepsilon}}. \tag{9}
\]

\noindent\textbf{Step 11 (Combine (8)–(9) into (7)).}  
Dividing both sides of (7) by $T$ yields the average regret bound
\[
\frac{1}{T}\sum_{k=0}^{T-1}\bigl(f(w_k)-f(w^\star)\bigr)
\le
\frac{D^{2}}{2\alpha}\frac{\sqrt{G_{T}+\varepsilon}}{T}
+\frac{\alpha G_{\max }^{2}}{2}\frac{1}{T}\sum_{k=0}^{T-1}\frac{1}{\sqrt{G_{k+1}+\varepsilon}}.
\tag{10}
\]
Equation \eqref{eq:update} together with the definition $G_{k+1}=G_k+g_k^{2}$ ensures $G_{k+1}\ge g_k^{2}$, so every denominator is at least $\sqrt{g_k^{2}}=|g_k|$. Consequently,
\[
\frac{1}{\sqrt{G_{k+1}+\varepsilon}}\le\frac{1}{|g_k|}\le\frac{1}{G_{\max }},
\]
and the second term in (10) is $O(\alpha G_{\max }/T)$.  
Choosing $\alpha = D/(G_{\max }\sqrt{T})$ balances the two terms, giving
\[
\frac{1}{T}\sum_{k=0}^{T-1}\bigl(f(w_k)-f(w^\star)\bigr)
\le
\frac{DG_{\max }}{\sqrt{T}} + O\!\left(\frac{1}{T}\right).
\]
This completes the proof. \hfill $\square$

%%%%%%%%%%%%%%%%%%%%%%%%%%%%%%%%%%%%%%%%%%%%%%%%%%%%%%%%%%%%%%%%%
\section*{Rate Derivation (With Constants)}
From the final bound,
\[
\mathbb{E}\bigl[f(\bar w_T)-f(w^\star)\bigr]
\le
\underbrace{\frac{DG_{\max }}{\sqrt{T}}}_{\text{dominant term}}
\;+\;
\underbrace{\frac{D^{2}\sqrt{\varepsilon}}{2\alpha T}}_{\text{numerical‑stability term}}
\;+\;
\underbrace{\frac{\alpha G_{\max }^{2}}{2T}\sum_{k=0}^{T-1}\frac{1}{\sqrt{G_{k+1}+\varepsilon}}}_{\text{gradient‑variance term}}.
\]
If $\varepsilon$ is chosen as $10^{-8}$, the second term is negligible for any reasonable $T$.  
The third term can be upper‑bounded by $\frac{\alpha G_{\max }}{2\sqrt{T}}$ using the monotonicity of $G_{k+1}$.  
Hence a compact constant‑free expression:
\[
\mathbb{E}\bigl[f(\bar w_T)-f(w^\star)\bigr]\le
\frac{DG_{\max }+\alpha G_{\max }/2}{\sqrt{T}}+o\!\left(\frac{1}{\sqrt{T}}\right).
\]

%%%%%%%%%%%%%%%%%%%%%%%%%%%%%%%%%%%%%%%%%%%%%%%%%%%%%%%%%%%%%%%%%
\section*{Special Cases}
\begin{enumerate}[label=\alph*.]
\item \textbf{Exact minimizer attained in finite steps.}  
If at some iteration $k_0$ we obtain $g_{k_0}=0$, then $G_{k}\equiv G_{k_0}$ for all $k\ge k_0$ and the update becomes $w_{k+1}=w_k$. The algorithm stops with $w_{k_0}=w^\star$ and the bound becomes zero.

\item \textbf{Strongly convex $f$.}  
If $f$ is $\mu$‑strongly convex, one can replace the average‑regret bound by a pointwise bound
\[
f(w_T)-f(w^\star)\le \frac{D^{2}}{2\alpha}\frac{\sqrt{G_T+\varepsilon}}{T},
\]
which after the same choice of $\alpha$ yields an $O(\log T /T)$ rate, matching the optimal rate for AdaGrad under strong convexity.

\item \textbf{Unbounded domain but bounded gradients.}  
If $\mathcal{W}=\R$ but $|g_k|\le G_{\max }$ holds, the same proof goes through with $D$ replaced by $\|w_0-w^\star\|$, i.e. the bound depends only on the initial distance.

\end{enumerate}

%%%%%%%%%%%%%%%%%%%%%%%%%%%%%%%%%%%%%%%%%%%%%%%%%%%%%%%%%%%%%%%%%
\section*{Conclusion}
We have presented a complete, step‑by‑step convergence analysis of a scalar AdaGrad‑style method that adapts its stepsize via the accumulated squared gradient.  Starting from the convexity inequality, we derived an exact regret bound, identified the optimal stepsize choice, and obtained the classic $O(1/\sqrt{T})$ convergence rate.  The proof is fully self‑contained, each algebraic manipulation is explicit, and all non‑trivial steps are annotated with line numbers for easy verification.

\end{document}